% 第 5 章：函数专题分层训练（MathNook 60 题候选题池，逐题验算后编入）
\chapter{函数专题分层训练}

\begin{center}
{\small\itshape 题号保留 M/A/K/P/B/T/H 六组稳定编号，便于分层布置；带“选做”者供学有余力者使用。}
\end{center}

\section{摸底诊断（M）}

\begin{problem}\level{易}
设 $A=[-1,3)$，$B=(1,5]$，求 $A\cap B$。
\end{problem}
\begin{solution}
$A\cap B=\{x\mid 1<x<3\}=(1,3)$。
\end{solution}

\begin{problem}\level{易}
判断下列对应能否成为集合 $A$ 上的函数：（1）$A=\{-1,0,1\}$，$y=\sqrt{x+1}$；（2）$A=[0,2]$，$y=\dfrac{1}{x-1}$。
\end{problem}
\begin{solution}
（1）三个元素代入均有意义且唯一，是函数；（2）$x=1\in A$ 处无定义，故不是 $A$ 上的函数。
\end{solution}

\begin{problem}\level{中}
求 $f(x)=\dfrac{\sqrt{(x-1)(4-x)}}{x-2}$ 的定义域。
\end{problem}
\begin{solution}
需 $(x-1)(4-x)\ge0$ 且 $x\ne2$，即 $1\le x\le4$ 且 $x\ne2$，定义域 $[1,2)\cup(2,4]$。
\end{solution}

\begin{problem}\level{易}
设 $f(x)=x^2-2x+3$，求 $f(-1)$ 与 $f(a+1)$。
\end{problem}
\begin{solution}
$f(-1)=1+2+3=6$；$f(a+1)=(a+1)^2-2(a+1)+3=a^2+2$。
\end{solution}

\begin{problem}\level{易}
已知 $f(x)=x+2\ (x<1)$，$f(x)=x^2\ (x\ge1)$，求 $f(0)$、$f(2)$。
\end{problem}
\begin{solution}
$f(0)=0+2=2$；$f(2)=2^2=4$。
\end{solution}

\begin{problem}\level{中}
已知 $f(x)$ 的定义域为 $[-2,3]$，求 $f(2x-1)$ 的定义域。
\end{problem}
\begin{solution}
$-2\le 2x-1\le3\Rightarrow -\tfrac12\le x\le2$，即 $\left[-\tfrac12,2\right]$。
\end{solution}

\begin{problem}\level{易}
解方程 $x^2-5x+6=0$。
\end{problem}
\begin{solution}
$(x-2)(x-3)=0$，故 $x=2$ 或 $x=3$。
\end{solution}

\begin{problem}\level{易}
已知 $f$ 在 $\mathbb{R}$ 上严格递增，$f(2)=5$，解不等式 $f(x)>5$。
\end{problem}
\begin{solution}
由严格递增，$f(x)>f(2)\iff x>2$，解集 $(2,+\infty)$。
\end{solution}

\section{基础计算专项（A）}

\begin{problem}\level{易}
化简 $(2x-3)(x+4)-(x-1)^2$。
\end{problem}
\begin{solution}
原式 $=2x^2+5x-12-(x^2-2x+1)=x^2+7x-13$。
\end{solution}

\begin{problem}\level{易}
因式分解 $x^2-7x+12$。
\end{problem}
\begin{solution}
$x^2-7x+12=(x-3)(x-4)$。
\end{solution}

\begin{problem}\level{易}
配方 $2x^2-8x+3$。
\end{problem}
\begin{solution}
$2x^2-8x+3=2(x-2)^2-5$。
\end{solution}

\begin{problem}\level{易}
化简 $\dfrac{x^2-9}{x-3}$ 并注明条件。
\end{problem}
\begin{solution}
$\dfrac{(x-3)(x+3)}{x-3}=x+3$，条件是 $x\ne3$。
\end{solution}

\begin{problem}\level{易}
化简 $\dfrac{1}{x-1}-\dfrac{1}{x+1}$ 并注明条件。
\end{problem}
\begin{solution}
$\dfrac{(x+1)-(x-1)}{(x-1)(x+1)}=\dfrac{2}{x^2-1}$，条件是 $x\ne\pm1$。
\end{solution}

\begin{problem}\level{中}
化简 $\dfrac{x^2-5x+6}{x-2}$ 并注明定义域。
\end{problem}
\begin{solution}
$\dfrac{(x-2)(x-3)}{x-2}=x-3$，定义域 $x\ne2$。
\end{solution}

\begin{problem}\level{易}
解不等式 $x^2-5x+6<0$。
\end{problem}
\begin{solution}
$(x-2)(x-3)<0\Rightarrow 2<x<3$。
\end{solution}

\begin{problem}\level{中}
解不等式 $\dfrac{x-2}{x+1}\ge0$。
\end{problem}
\begin{solution}
分子分母同号且分母非零，得 $x<-1$ 或 $x\ge2$。
\end{solution}

\begin{problem}\level{中}
解不等式 $|2x-3|\le5$。
\end{problem}
\begin{solution}
$-5\le 2x-3\le5\Rightarrow -1\le x\le4$。
\end{solution}

\begin{problem}\level{易}
求 $y=\sqrt{2x+6}$ 的定义域。
\end{problem}
\begin{solution}
$2x+6\ge0\Rightarrow x\ge-3$。
\end{solution}

\begin{problem}\level{易}
设 $f(x)=x^2-3x+1$，求 $f(x+1)$。
\end{problem}
\begin{solution}
$f(x+1)=(x+1)^2-3(x+1)+1=x^2-x-1$。
\end{solution}

\begin{problem}\level{中}
判断 $x^2-4x+1\ge-3$ 是否对任意实数成立。
\end{problem}
\begin{solution}
等价于 $x^2-4x+4=(x-2)^2\ge0$，恒成立。
\end{solution}

\section{函数核心例题（K）}

\begin{problem}\level{中}
判断 $f(x)=\dfrac{x^2-1}{x-1}$ 与 $g(x)=x+1$ 是否为同一个函数。
\end{problem}
\begin{solution}
不是。$f$ 的定义域为 $x\ne1$，$g$ 的定义域为 $\mathbb{R}$，定义域不同。
\end{solution}

\begin{problem}\level{中}
已知 $f(x)$ 的定义域为 $[-1,2]$，求 $f(x^2)$ 的定义域。
\end{problem}
\begin{solution}
$-1\le x^2\le2\Rightarrow x^2\le2\Rightarrow -\sqrt2\le x\le\sqrt2$。
\end{solution}

\begin{problem}\level{中}
设 $f(x)=\sqrt{4-x}$，求 $f(x^2-1)$ 的解析式与定义域。
\end{problem}
\begin{solution}
$f(x^2-1)=\sqrt{5-x^2}$；由 $5-x^2\ge0$ 得 $-\sqrt5\le x\le\sqrt5$。
\end{solution}

\begin{problem}\level{难}
已知 $f(3-2x)$ 的定义域为 $[0,2]$，求 $f(x)$ 的定义域，并说明能否唯一确定 $f$。
\end{problem}
\begin{solution}
$x\in[0,2]\Rightarrow 3-2x\in[-1,3]$，故 $f$ 的定义域为 $[-1,3]$；但 $f$ 的对应法则无法由此唯一确定。
\end{solution}

\begin{problem}\level{中}
设 $f(x)=x^2+4|x|+1$。（1）改写成含 $(|x|+2)^2$ 的形式；（2）判断奇偶性；（3）比较 $f(2)$ 与 $f(1)$。
\end{problem}
\begin{solution}
（1）$f(x)=(|x|+2)^2-3$；（2）只含 $|x|$，为偶函数；（3）$f(2)=13>f(1)=6$。
\end{solution}

\begin{problem}\level{中}
求 $x^2-4x+3$ 在 $[0,3]$ 上的最大值与最小值。
\end{problem}
\begin{solution}
$x^2-4x+3=(x-2)^2-1$，最小值 $-1$（$x=2$）；最大值 $3$（$x=0$）。
\end{solution}

\begin{problem}\level{易}
把 $y=x^2$ 向右平移 $2$ 个单位、向上平移 $1$ 个单位，求新解析式与顶点。
\end{problem}
\begin{solution}
$y=(x-2)^2+1$，顶点 $(2,1)$。
\end{solution}

\begin{problem}\level{中}
用定义证明 $f(x)=3x-1$ 在 $\mathbb{R}$ 上严格递增。
\end{problem}
\begin{solution}
任取 $x_1<x_2$，$f(x_2)-f(x_1)=3(x_2-x_1)>0$，故严格递增。
\end{solution}

\section{上午知识复测（P）}

\begin{problem}\level{中}
求 $y=\dfrac{\sqrt{(x+1)(3-x)}}{x-1}$ 的定义域。
\end{problem}
\begin{solution}
$(x+1)(3-x)\ge0$ 且 $x\ne1$，得 $[-1,1)\cup(1,3]$。
\end{solution}

\begin{problem}\level{易}
化简 $\dfrac{x^2-4}{x+2}$ 并注明条件。
\end{problem}
\begin{solution}
$=x-2$，条件是 $x\ne-2$。
\end{solution}

\begin{problem}\level{中}
若 $(x-a)^2$ 在 $[0,2]$ 上最小值为 $0$，求 $a$ 的范围。
\end{problem}
\begin{solution}
需 $a\in[0,2]$。
\end{solution}

\begin{problem}\level{中}
已知 $f$ 在 $\mathbb{R}$ 上严格递增，解 $f(2x+1)>f(5-x)$。
\end{problem}
\begin{solution}
$2x+1>5-x\Rightarrow x>\dfrac43$。
\end{solution}

\begin{problem}\level{易}
把 $x^2+2|x|+1$ 写成平方形式并判断奇偶性。
\end{problem}
\begin{solution}
$x^2+2|x|+1=(|x|+1)^2$，是偶函数。
\end{solution}

\section{综合计算专项（B）}

\begin{problem}\level{中}
分解因式 $(2x-1)^2-(x+2)^2$。
\end{problem}
\begin{solution}
$=[(2x-1)-(x+2)][(2x-1)+(x+2)]=(x-3)(3x+1)$。
\end{solution}

\begin{problem}\level{易}
化简 $\dfrac{x^2-4x+3}{x-1}$ 并注明条件。
\end{problem}
\begin{solution}
$=\dfrac{(x-1)(x-3)}{x-1}=x-3$，条件是 $x\ne1$。
\end{solution}

\begin{problem}\level{易}
化简 $\dfrac{x^2-a^2}{x-a}$ 并注明条件。
\end{problem}
\begin{solution}
$=x+a$，条件是 $x\ne a$。
\end{solution}

\begin{problem}\level{中}
化简 $\dfrac{2}{x-1}-\dfrac{1}{x+1}$。
\end{problem}
\begin{solution}
$=\dfrac{2(x+1)-(x-1)}{(x-1)(x+1)}=\dfrac{x+3}{x^2-1}$，$x\ne\pm1$。
\end{solution}

\begin{problem}\level{中}
配方 $x^2-2(a+1)x+(a+1)^2+2$。
\end{problem}
\begin{solution}
$=[x-(a+1)]^2+2$。
\end{solution}

\begin{problem}\level{中}
求 $x^2-6x+5$ 在 $[1,4]$ 上的最大值与最小值。
\end{problem}
\begin{solution}
$x^2-6x+5=(x-3)^2-4$，最小值 $-4$（$x=3$），最大值 $0$（$x=1$）。
\end{solution}

\begin{problem}\level{中}
若 $(x-a)^2$ 在 $[-1,3]$ 上最小值为 $0$，求 $a$ 的范围。
\end{problem}
\begin{solution}
$a\in[-1,3]$。
\end{solution}

\begin{problem}\level{中}
把 $|x-2|+|x+1|$ 写成分段函数。
\end{problem}
\begin{solution}
$\begin{cases}-2x+1,&x<-1,\\ 3,&-1\le x\le2,\\ 2x-1,&x>2.\end{cases}$
\end{solution}

\begin{problem}\level{中}
解不等式 $\dfrac{x-3}{x+2}<0$。
\end{problem}
\begin{solution}
$-2<x<3$。
\end{solution}

\begin{problem}\level{难}
解不等式 $(|2x-1|+1)^2>(|x+2|+1)^2$，并说明去平方的合法性。
\end{problem}
\begin{solution}
两边均为正，平方保持大小关系，等价于 $|2x-1|>|x+2|$，再平方得 $(x-3)(3x+1)>0$，故 $x<-\tfrac13$ 或 $x>3$。
\end{solution}

\begin{problem}\level{中}
把 $f(x)=x-\dfrac{2}{x-1}$ 写成以 $x-1$ 为整体的形式。
\end{problem}
\begin{solution}
$f(x)=(x-1)-\dfrac{2}{x-1}+1$。
\end{solution}

\begin{problem}\level{难}
求 $q_a(x)=(x-a)^2+1$ 在 $[0,4]$ 上的最大值，按参数 $a$ 分类。
\end{problem}
\begin{solution}
最远端点决定最大值：$a\le2$ 时取 $x=4$，最大 $(4-a)^2+1$；$a\ge2$ 时取 $x=0$，最大 $a^2+1$。
\end{solution}

\begin{problem}\level{中}
设 $f(x)=\dfrac{x-1}{x+2}$，对 $-2<x_1<x_2$ 求 $f(x_2)-f(x_1)$ 并判断符号。
\end{problem}
\begin{solution}
$f(x)=1-\dfrac{3}{x+2}$，差 $=3\left(\dfrac{1}{x_1+2}-\dfrac{1}{x_2+2}\right)>0$，为正。
\end{solution}

\begin{problem}\level{易}
（已学指数运算者选做）化简 $\dfrac{2^{x+1}}{2^{x-2}}$。
\end{problem}
\begin{solution}
同底相除指数相减：$2^{(x+1)-(x-2)}=2^3=8$。
\end{solution}

\section{性质与核心例题（T）}

\begin{problem}\level{中}
已知 $f$ 在 $\mathbb{R}$ 上严格递增，证明 $f(x)>f(a)\iff x>a$（充分性与必要性分别证）。
\end{problem}
\begin{solution}
充分性：$x>a\Rightarrow f(x)>f(a)$；必要性：若 $x\le a$ 则 $f(x)\le f(a)$，矛盾。
\end{solution}

\begin{problem}\level{难}
已知 $f$ 在 $\mathbb{R}$ 上严格递增，解 $f(x^2-3x)>f(2x-4)$。
\end{problem}
\begin{solution}
等价于 $x^2-3x>2x-4$，即 $(x-1)(x-4)>0$，故 $x<1$ 或 $x>4$。
\end{solution}

\begin{problem}\level{难}
设 $f(x)=x(x+4)+1\ (x\ge0)$，$f(x)=x(x-4)+1\ (x<0)$。（1）用绝对值统一；（2）判断奇偶性；（3）解 $f(2a)>f(a-3)$。
\end{problem}
\begin{solution}
（1）$f(x)=x^2+4|x|+1$；（2）偶函数；（3）$f$ 在 $[0,+\infty)$ 递增，故等价于 $|2a|>|a-3|$，即 $(a+3)(3a-3)>0$，解得 $a<-3$ 或 $a>1$。
\end{solution}

\begin{problem}\level{难}
不用导数证明 $f(x)=x-\dfrac{1}{x-1}$ 在 $(1,+\infty)$ 上严格递增。
\end{problem}
\begin{solution}
设 $1<x_1<x_2$，差 $=(x_2-x_1)+\left(\dfrac{1}{x_1-1}-\dfrac{1}{x_2-1}\right)=(x_2-x_1)\left[1+\dfrac{1}{(x_1-1)(x_2-1)}\right]>0$。
\end{solution}

\begin{problem}\level{难}
设 $q_a(x)=(x-a)^2+1$，求它在 $[-1,3]$ 上的最大值与最小值，按 $a$ 分类。
\end{problem}
\begin{solution}
最小值：$a<-1$ 时 $(a+1)^2+1$；$-1\le a\le3$ 时 $1$；$a>3$ 时 $(a-3)^2+1$。
最大值：$a<-1$ 时 $(3-a)^2+1$；$-1\le a<1$ 时 $(3-a)^2+1$；$1\le a\le3$ 时 $(a+1)^2+1$；$a>3$ 时 $(a+1)^2+1$。
\end{solution}

\begin{problem}\level{难}
设 $f(x)=\dfrac{ax-1}{x+2}$。（1）写成 $a+\dfrac{k}{x+2}$；（2）讨论其在 $(-1,+\infty)$ 上的单调性；（3）讨论 $a=-\dfrac12$ 的情形。
\end{problem}
\begin{solution}
（1）$f(x)=a-\dfrac{2a+1}{x+2}$；（2）在 $(-1,+\infty)$ 上 $x+2$ 递增且为正，故 $a>-\tfrac12$ 时 $f$ 递增，$a<-\tfrac12$ 时递减；（3）$a=-\tfrac12$ 时 $f\equiv-\dfrac12$ 为常函数。
\end{solution}

\begin{problem}\level{中}
设 $g(x)=(x-2)^2-1$，求顶点、对称轴、零点，并解 $g(x)\le0$。
\end{problem}
\begin{solution}
顶点 $(2,-1)$，对称轴 $x=2$，零点 $1,3$；$g(x)\le0\iff 1\le x\le3$。
\end{solution}

\begin{problem}\level{中}
设 $f$ 为偶函数且 $f(2)=3$。（1）求 $f(-2)$；（2）证明定义域关于原点对称；（3）若 $1$ 属于定义域，$-1$ 是否一定属于定义域？
\end{problem}
\begin{solution}
（1）$f(-2)=f(2)=3$；（2）偶函数要求对定义域内每个 $x$，$-x$ 也在定义域内；（3）一定。
\end{solution}

\section{综合提高（H）}

\begin{problem}\level{难}
求 $a$ 的范围，使得对任意 $x\in[-1,2]$，都存在 $y\in[0,9]$ 满足 $y=(x-a)^2$。
\end{problem}
\begin{solution}
等价于 $(x-a)^2\le9$ 对一切 $x\in[-1,2]$ 成立，即 $|x-a|\le3$，需 $a\in[x-3,x+3]$ 对所有 $x$ 成立，取交得 $a\in[-1,2]$。
\end{solution}

\begin{problem}\level{难}
求 $a$ 的范围，使得存在 $x\in[-1,2]$，对任意 $y\in[0,1]$ 都有 $(x-a)^2\ge y$；并与上一题比较量词顺序带来的差异。
\end{problem}
\begin{solution}
等价于存在 $x\in[-1,2]$ 使 $|x-a|\ge1$。由于 $[-1,2]$ 长 $3$ 而“到 $a$ 距离小于 $1$ 的点”至多覆盖长度 $2$，这样的 $x$ 总存在，故 $a\in\mathbb{R}$。与 H01 对照：$\forall x\,\exists y$ 对 $a$ 施加约束，$\exists x\,\forall y$ 反而无约束——量词顺序决定结论方向。
\end{solution}

\begin{problem}\level{难}
若 $f(x)=(x-a)^2+1$ 在 $[0,3]$ 上最小值为 $2$，求 $a$。
\end{problem}
\begin{solution}
需 $(x-a)^2$ 的最小值为 $1$。$a\in[0,3]$ 时最小值为 $0$，不合；$a<0$ 时最小在 $x=0$，$a^2=1\Rightarrow a=-1$；$a>3$ 时最小在 $x=3$，$(3-a)^2=1\Rightarrow a=4$。故 $a=-1$ 或 $a=4$。
\end{solution}

\begin{problem}\level{难}
求 $a$ 的范围，使得对任意 $x\in[1,2]$，都存在 $y\in[0,4]$ 满足 $y=x^2+a$。
\end{problem}
\begin{solution}
等价于 $x^2+a\in[0,4]$ 对一切 $x\in[1,2]$ 成立：上界要求 $a\le 4-x^2$（最严 $x=2$ 得 $a\le0$），下界要求 $a\ge -x^2$（最严 $x=1$ 得 $a\ge-1$）。故 $a\in[-1,0]$。
\end{solution}

\begin{problem}\level{难}
（提高选做）设 $f(x)=x-\dfrac{2}{x-1}$，求 $h,k$ 使 $g(x)=f(x+h)-k$ 为奇函数，并验证定义域的对称性。
\end{problem}
\begin{solution}
由 $f(x)=(x-1)-\dfrac{2}{x-1}+1$，令 $h=1$ 得 $f(x+1)=x-\dfrac{2}{x}+1$，再取 $k=1$ 得 $g(x)=x-\dfrac{2}{x}$，定义域 $x\ne0$ 关于原点对称且 $g(-x)=-g(x)$，确为奇函数。
\end{solution}
